\section{Механизм пропусков}
\label{sec:mnar}

Полные кривые снимали не для всех: отбор шёл по результату первичного скрининга. Это отбор по зависимой
переменной, то есть missing not at random по построению, и он же --- причина провала на CYP2D6,
измеренная в разделе~\ref{sec:d6}.

Спасает вот что: \emph{переменная отбора наблюдаема}. Скрининг есть для 4375 из 4905 соединений
обучающей выборки, и именно он определял, кого поднимут на подробное измерение.

\subsection{Как раскладывается совместное распределение}

Пишем совместную плотность через скрытую активность $a$, честно:
\begin{equation}
p(y_{\text{крив}}, y_{\text{скр}}, R \mid x,\theta)
= p(R \mid y_{\text{скр}}) \int p(a \mid x,\theta)\, p(y_{\text{скр}} \mid a)\, p(y_{\text{крив}} \mid a)^{R}\, da .
\label{eq:joint}
\end{equation}

Здесь $R$ --- индикатор попадания в подробное измерение, $p(a\mid x,\theta)$ --- что модель думает
об активности до измерений (центр --- предсказание сети, ширина --- признание того, что структура
определяет активность не полностью), два следующих множителя --- шум каждого прибора, зависящий только
от истинной активности. Множитель $p(y_{\text{крив}}\mid a)$ входит в степени $R$, то есть только для
отобранных.

Главный шаг --- вынос $p(R\mid y_{\text{скр}})$ за интеграл. Он законен потому, что решение об отборе
принималось глядя на показание скрининга, а не на скрытую активность: под интегралом по $a$ этот
множитель величины $a$ не содержит.

Применяя цепное правило, получаем рабочую форму:
\begin{equation}
p(y_{\text{крив}}, y_{\text{скр}}, R \mid \theta)
= \underbrace{p(R \mid y_{\text{скр}})}_{\text{без }\theta}
\cdot\; \underbrace{p(y_{\text{крив}} \mid y_{\text{скр}},\, \theta)}_{\text{условно}}
\cdot\; p(y_{\text{скр}} \mid \theta) .
\label{eq:factor}
\end{equation}

Первый множитель не зависит от $\theta$, при дифференцировании логарифма даёт ноль и полностью выпадает
из градиента. Значит можно обучаться на полном правдоподобии, вообще не моделируя процедуру отбора,
и оценки останутся состоятельными. Условие ровно одно: скрининг должен входить в модель как наблюдение.
Если его выбросить, первый множитель перестанет быть отделимым и смещение вернётся. По сути это поправка
хекмановского типа, у которой первая ступень не моделируется, а наблюдается.

Средний множитель обязан быть \emph{условным} по $y_{\text{скр}}$, и это место, где легко ошибиться.
Заменить его на маргинальный $p(y_{\text{крив}}\mid\theta)$ можно только при нулевом собственном латентном
разбросе соединения: два показания меряют одну скрытую активность, они коррелированы, а отбор идёт
по одному из них --- значит среди отобранных распределение второго сдвинуто. Маргинальный множитель
про этот сдвиг не знает.

\subsection{Что даёт симуляция}

Проверено на данных, где истина известна по построению: $a = x^{\top}\beta + \varepsilon$, скрининг у всех,
кривая --- у верхних двадцати восьми процентов по скринингу (порог квантиля 0.72).

\begin{table}[htbp]\centering\small
\begin{tabularx}{\textwidth}{@{}Xrr@{}}
\toprule
Способ подгонки & Смещение, норма & Наклон калибровки \\
\midrule
Только отобранные метки кривых & 0.433 & 0.731 \\
Совместное правдоподобие, отбор по скринингу & 0.019 & 1.012 \\
Только скрининг, все строки (контроль сверху) & 0.002 & 1.000 \\
Отбор по скрытой истине, скрининг у неотобранных не виден & 0.512 & 0.681 \\
\bottomrule
\end{tabularx}
\caption{Сорок повторов, шесть тысяч соединений, в подробное измерение попадает 28 процентов. Наклон
единица означает, что коэффициенты восстановлены без систематического сжатия.}
\label{tab:mnar}
\end{table}

Подгонка по одним отобранным меткам сжимает коэффициенты до 0.73 от истины --- ровно то ограничение
диапазона, которое видно на CYP2D6. Совместное правдоподобие возвращает 1.01.

Четвёртая строка важнее первых трёх. Сломайте условие, то есть сделайте отбор зависящим от скрытой
величины и уберите скрининг у неотобранных, и совместная схема поломается так же, как наивная. Значит
аргумент держится не на общей идее «использовать больше данных», а на наблюдаемости первой ступени
отбора. В наших данных она наблюдаема.

Отдельно измерена цена замены условного множителя на маргинальный. Смещения нет при нулевом латентном
остатке, при $\sigma_{\text{лат}} = 0.35$ оно примерно втрое меньше стандартной ошибки оценки, при
$\sigma_{\text{лат}} = 0.70$ уже сравнимо с ней. И независимо от смещения такая схема занижает собственную
погрешность в 1.18--1.21 раза, потому что считает два коррелированных наблюдения одного соединения
за два независимых свидетельства. Значит голова, предсказывающая $\pic$, должна получать на вход показание скрининга того же соединения.
Или, что то же самое, обе величины должны выводиться из общего латентного слоя с явной ковариацией,
а не из двух независимых голов.

На инференсе скрининга нет, но это не мешает. Условный множитель нужен на обучении, чтобы оценка
не была смещена отбором; на предсказании мы маргинализуем по $y_{\text{скр}}$, что при общем латентном
слое сводится к одному прогону сети.

Оговорка: 530 соединений с кривыми не имеют скрининга, и для них аргумент не работает. По выборке
целиком доля небольшая, но лежит она не равномерно: все 530 --- это метки 3A4, почти четверть всех
кривых этого фермента. Так что на трёх ферментах оговорка пустая, а на 3A4 её надо закрывать --- взвесить
или выкинуть в абляции.
